\textbf{Pour cette première partie, aucune justification n'est demandée et une seule réponse est possible par question.} Les justifications ci-dessous sont fournies à titre pédagogique.

Une réponse fausse ou l'absence de réponse n'enlève aucun point : il ne faut donc jamais laisser une question sans réponse.

\subsection*{Question 1 \quad Réponse B}

On applique la règle $a^m \times a^n = a^{m+n}$ :
\[3^4 \times 3^2 = 3^{4+2} = 3^6.\]

\subsection*{Question 2 \quad Réponse B}

Le coefficient directeur d'une droite passant par deux points A$(x_A ; y_A)$ et
B$(x_B ; y_B)$ est
\[a = \dfrac{y_B - y_A}{x_B - x_A} = \dfrac{6-2}{3-1} = \dfrac{4}{2} = 2.\]

\subsection*{Question 3 \quad Réponse C}

Par lecture graphique, on repère les points dont l'ordonnée vaut $14,5$ : ce sont ceux
d'abscisses 2018 et 2021.

\subsection*{Question 4 \quad Réponse C}

On applique l'identité remarquable $(a-b)^2 = a^2 - 2ab + b^2$ avec $a = x$ et $b = 5$ :
\[(x-5)^2 = x^2 - 2\times 5 \times x + 5^2 = x^2 - 10x + 25.\]

\subsection*{Question 5 \quad Réponse C}

La moyenne est la somme des notes divisée par leur nombre :
\[\overline{x} = \dfrac{14+11+11+4+10}{5} = \dfrac{50}{5} = 10.\]

\subsection*{Question 6 \quad Réponse C}

La droite coupe l'axe des ordonnées en 3 : l'ordonnée à l'origine est $b = 3$. Elle passe par les point de coordonnées $(0 ; 3)$ et $(-2 ; 0)$, donc son coefficient directeur vaut
\[a = \dfrac{3 - 0}{0 - (-2)} = \dfrac{3}{2}.\]
D'où $f(x) = \dfrac{3}{2}x + 3$.

\subsection*{Question 7 \quad Réponse C}

$-10x + 4 = 0 \iff x = \dfrac{4}{10} = 0,4$. Le coefficient de $x$ vaut $-10 < 0$ : la
fonction affine est décroissante, donc positive avant $0,4$ et négative après.

\subsection*{Question 8 \quad Réponse A}

Le taux de réduction se calcule par rapport au prix initial :
\[\dfrac{20-25}{25} = \dfrac{-5}{25} = -\dfrac{1}{5} = -0,20 = -20\,\%.\]
Le pourcentage de réduction est de $20\,\%$.
